\title{Large Language Models Can Automatically Engineer Features for Few-Shot Tabular Learning}

\begin{abstract}
Large Language Models (LLMs) have demonstrated exceptional proficiency in addressing complex and novel reasoning problems, revealing significant promise for tabular learning—a critical component for numerous practical tasks. This paper introduces \textsf{FeatLLM}, a new in-context learning paradigm where LLMs operate as automated feature engineers to create transformed input datasets optimized for subsequent tabular prediction. The new features extracted by the LLM are then utilized by a straightforward machine learning method, such as linear regression, for predicting class probabilities, thereby achieving robust few-shot learning results. Notably, \textsf{FeatLLM} restricts inference-time computations to the use of this simple predictive model leveraging the engineered features, circumventing the necessity of consulting the LLM for each prediction. In contrast to other LLM-driven strategies, \textsf{FeatLLM} avoids inference-time queries to the LLM, depends solely on API-based model access, and is not hindered by prompt size restrictions. Empirical evidence across a broad array of tabular datasets, sourced from various domains, illustrates that \textsf{FeatLLM} systematically crafts superior features—yielding an average improvement of 10\% over contemporaries such as TabLLM and STUNT.
\end{abstract}

\section{Introduction}
The advent of Large Language Models (LLMs) has seen substantial breakthroughs in their capability to generate appropriate solutions for previously unseen challenges, all without the need for training specialized to a given task~\cite{creswell2022selection,imani2023mathprompter,taylor2022galactica}. These models, benefiting from vast and diverse text datasets, encapsulate a wealth of prior knowledge, frequently replacing the requirement for specialized human expertise in various sectors~\cite{Genomes2021,singhal2023large,wilcox2003role}. Such qualities make LLMs particularly valuable for streamlining tasks within a wide spectrum of applications, ranging from dialogue evaluation~\cite{finch2023leveraging} to the automated repair of source code~\cite{fan2023automated}. Especially noteworthy is that, by leveraging task outlines and providing minimal examples within the prompts, LLMs can capture task context, recognizing salient features and demonstrating nuanced attention to relevant data points~\cite{brown2020language,zhao2023survey}. This capacity highlights the potential for LLMs to function adeptly as automated feature engineering tools.

Recently, researchers have expanded the application of LLMs' inferential and knowledge-based advantages into tabular data settings. For example, an LLM’s understanding that certain health conditions are more probable in older adults can aid predictive modeling in healthcare. Contemporary approaches formulate task- and feature-related instructions in plain language, serialize tabular inputs, and subsequently process this information via an LLM~\cite{dinh2022lift,wang2023anypredict}. Following this, frameworks typically employ either in-context learning methodologies~\cite{nam2023semi} or make use of efficient fine-tuning strategies~\cite{dinh2022lift,hegselmann2023tabllm}. These efforts illustrate that the injection of LLM-driven knowledge confers advantages that consistently surpass standard tabular learning baselines, particularly in data-constrained or few-shot scenarios~\cite{hegselmann2023tabllm}.

Contemporary methods leveraging LLMs for tabular data face notable shortcomings. In end-to-end prediction settings, these approaches necessitate conducting at least one LLM inference for each individual instance, which incurs substantial computational overhead. The majority demand fine-tuning of the LLM to attain strong predictive performance, which is impractical for LLMs that lack public access to their training infrastructure or tuning capabilities. This constraint has become increasingly significant, as many high-performing, recently introduced LLMs restrict users to inference-only APIs~\cite{anil2023palm,openai2023gpt4}. Additionally, many of these strategies struggle with prompt scalability~\cite{liu2023lost}; as the feature dimensionality of tabular data increases and serialized prompts grow longer, these methods often become intractable or exhibit pronounced drops in effectiveness. Collectively, these obstacles complicate the deployment of LLM-based tabular systems in practical, real-world scenarios.

To address these issues, our method reframes the function of LLMs, shifting from direct end-to-end prediction to the task of feature engineering, thereby facilitating more effective use of pretrained knowledge. We present \textsf{FeatLLM}, a new in-context learning scheme, which is designed to uncover the implicit ``criteria'' LLMs employ in their decision-making processes. For example, when applied to medical diagnosis, the LLM can infer and articulate explicit rules—linking specific feature thresholds or combinations to disease identification. These distilled rules become new, meaningful features, replacing raw ones for downstream predictive tasks. By relocating the LLM's role to feature construction, a downstream pipeline can be streamlined, reducing dependence on resource-intensive models during inference and enhancing response times. Utilizing the in-context learning capabilities of LLMs, we demonstrate that such features can be extracted simply by incorporating example demonstrations into prompts, circumventing any further model training. Finally, by integrating ensemble learning techniques such as feature bagging~\cite{breiman2001random}, our approach can control prompt length, helping to overcome challenges related to lengthy serialized inputs.

\textsf{FeatLLM} decomposes the process of creating new features from prompts into two main sub-tasks: (1) comprehending the core problem and analyzing how features relate to target variables; and (2) synthesizing clear and discriminative rules for class separation, drawing on both prior insight and limited example instances. This methodical reasoning framework prompts LLMs to disregard features that lack relevance during rule formulation. The identified rules are subsequently applied to the dataset—where data instances are processed as code segments—resulting in binary indicators that reflect whether each sample fulfills the corresponding rule. Ultimately, a simple predictive model is trained on these binary features to estimate class membership probabilities. Overall system robustness is augmented by leveraging ensemble strategies: feature engineering is conducted multiple times, with random subsets of features and samples selected via bagging. By controlling the amount of features and data involved during this step, \textsf{FeatLLM} curtails potential issues stemming from overly large prompts. Operating entirely without the need for further training, and incurring minimal inference costs, this approach enables LLMs to address a broad spectrum of practical challenges.

The effectiveness of \textsf{FeatLLM} is assessed across 13 tabular datasets under scenarios with limited data. Experimental results highlight its ability to extract informative features by integrating background knowledge with empirical insights. Our approach consistently achieves superior performance compared to state-of-the-art few-shot learning baselines over diverse experimental configurations. The implementation is publicly accessible via an anonymized GitHub repository at~\texttt{https://github.com/Sungwon-Han/FeatLLM}.

\section{Related Work}

\subsection{Few-Shot Learning with Tabular Data} The field of few-shot learning has made it possible to derive representations from data that generalize well, even when labeling resources are scarce~\cite{chen2018closer}. Although originally developed in the context of computer vision, few-shot learning approaches have demonstrated strong generalization ability across multiple data types, such as text, audio, and, more recently, tabular datasets~\cite{majumder2022few,nam2023stunt,schick2021exploiting}. However, working with very limited labeled data introduces a heightened danger of picking up non-causal or coincidental associations~\cite{bartlett2021deep}. To address this, recent research has turned to auxiliary information sources and the explicit encoding of domain knowledge to guide models with suitable inductive biases during training. Strategies in this direction include leveraging abundant unlabeled data in combination with data augmentation within a semi-supervised framework~\cite{ucar2021subtab,yoon2020vime} and adopting unsupervised, task-independent pre-training methods to initialize models effectively~\cite{somepalli2021saint}. Commonly, such pre-training schemes rely on objectives like masking or intentional corruption of input features, followed by reconstructing the masked elements~\cite{arik2021tabnet,majmundar2022met}, or by employing data augmentation alongside contrastive learning losses~\cite{bahri2022scarf,chen2020simple}. Despite these advances, the structural diversity and heterogeneity in tabular datasets pose significant obstacles for defining universally effective augmentations, limiting the success of these techniques in the few-shot regime~\cite{nam2023stunt}. 

More recent research has introduced models trained across a wide array of tabular datasets, which need not be directly related to the ultimate task, enabling knowledge transfer to target applications~\cite{zhang2023generative,levin2022transfer,zhu2023xtab}. TransTab~\cite{wang2022transtab}, for example, utilizes transformer architectures to extract semantic links between table columns by using column names as additional context, beyond mere cell values. The generalizable patterns learned from assorted table-column structures enable both zero-shot and few-shot performance on unseen tasks. Similarly, TabPFN~\cite{hollmann2022tabpfn} generates synthetic data from an assortment of probability distributions, utilizing this data to pre-train a Bayesian neural network. At test time, TabPFN accepts both support (labeled) and query (unlabeled) data from the target task and performs inference directly, without fine-tuning. The multi-source prior established during pre-training enables the method to outperform conventional tree-based algorithms, particularly for low-data scenarios.

\subsection{Language-Interfaced Tabular Learning} Integrating large language models (LLMs)~\cite{anil2023palm,openai2023gpt4} presents another compelling opportunity for infusing prior knowledge into tabular learning. LLMs are pre-trained on vast quantities of textual data and have illustrated robust performance in transfer to previously unseen domains~\cite{brown2020language}. Recent work has begun exploring the application of such models for tabular tasks by representing structured data in textual form, which is then embedded into LLM prompts~\cite{dinh2022lift,hegselmann2023tabllm,wang2023anypredict}. Input prompts constructed from tabular entries, feature definitions, and problem statements allow LLMs to reason over both relational structure and task requirements. Methodologies in this space have included fine-tuning with parameter-efficient adaptations like LoRA~\cite{hu2021lora} or IA3~\cite{liu2022few}, as well as using in-context learning, where several annotated examples are provided in the prompt to enable inference by the pretrained model, without further weight updates~\cite{dinh2022lift,hegselmann2023tabllm,nam2023semi,slack2023tablet}.

While prior approaches have demonstrated success in enhancing few-shot learning on tabular data, their dependency on LLMs for inference in every prediction limits practical adoption in real-world industrial contexts. This research moves away from the typical reliance on LLMs as opaque, end-to-end prediction tools for individual samples. Rather, our focus is on utilizing LLMs to explicitly infer the underlying conditions or rules defining each class. In our approach, \textsf{FeatLLM} interprets these inferred rules into programmatic logic to automatically engineer new features, thereby removing the necessity to route each prediction through the LLM. This leverages the advantages of in-context learning: extracting and synthesizing rules by using both existing knowledge and information from the limited provided examples.

\section{Methods}
The central function of \textsf{FeatLLM} is to identify and extract “rules”—defining criteria that differentiate each answer class—based on the limited samples presented to it, as illustrated in Figure~\ref{fig:main_model}. The LLM is tasked with formulating these rules, which are then parsed and converted into code that generates new binary indicator features for each sample, signifying whether each condition holds. These constructed binary features serve as input for a weighted dot product operation, with the weights constrained to be non-negative and optimized during model training, to yield class probability estimates. Through the training process, we determine both the relevance of each engineered feature and its impact on predicting class membership (refer to Section~\ref{Sec:training}). This methodology is iteratively applied using bagging for variance reduction, with the final predictions obtained through an ensemble of such models. The problem setup and comprehensive procedural details for each stage are elaborated in the following sections.

\vspace*{-0.12in}
\paragraph{Problem formulation.} We focus on a tabular dataset comprising $N$ labeled entries, denoted as $\mathcal{D} = \{(\mathbf{x}^i, \mathbf{y}^i)\}_{i=1}^{N}$. In this setup, each data point is characterized by $d$ features, with every $\mathbf{x}^i$ represented as a vector of $d$ dimensions. These features are associated with natural language names, specified by the set $F = \{f_j\}_{j=1}^{d}$, and their accompanying definitions are specified elsewhere. In the context of classification, the target variable $\mathbf{y}^i$ takes a value from a set of specified classes. For $k$-shot learning scenarios, the model is trained on only $k$ ($<N$) labeled examples that are randomly drawn from $\mathcal{D}$.

\subsection{Prompt Design for Extracting Rules}
\label{Sec:prompt_design}
To improve the accuracy of rule extraction using LLMs, we structure the reasoning steps in a way that mirrors expert approaches to tabular learning. The prompt is structured into three key segments (visualized in Fig.~\ref{fig:prompt_for_rule} and detailed in Appendix~\ref{sec:example_rule}):

\vspace*{-0.12in}
\paragraph{Basic information description.} This segment lays out the fundamental context needed to approach the task (see orange highlighting in Fig.~\ref{fig:prompt_for_rule}). It outlines the task, incorporates label information, provides feature descriptions, and supplies illustrative demonstrations. The task itself is introduced as a specific question (for example, ``Does this patient’s myocardial infarction complications data indicate chronic heart failure? Yes or no?"). Supplementary instances can be found in Appendix~\ref{sec:task_desc}. Each feature’s description specifies its data type and, optionally, its definition. Categorical features can be accompanied by sample categories. For demonstrations, a small number of labeled training cases are transformed into text following a serialization process. This serialization mirrors prior methods~\cite{dinh2022lift,hegselmann2023tabllm}:

\begin{align}
&\text{Serialize}(\mathbf{x}^i,\mathbf{y}^i, F) = \nonumber \\ 
&``\ f_1\ \text{is}\ \mathbf{x}^i_1.\ ... \ f_d\  \text{is}\  \mathbf{x}^i_d.\ \ \text{Answer:}\  \mathbf{y}^i\ ”
\end{align}

\vspace*{-0.12in}
\paragraph{Reasoning instruction.} Our approach seeks to strengthen the reasoning capabilities of the LLM by offering explicit reasoning instructions (refer to the blue annotations in Fig.~\ref{fig:prompt_for_rule}). This guidance starts with an introductory sentence modeled after the chain-of-thought paradigm~\cite{wei2022chain}. Subsequently, we delineate the LLM's inference process into two distinct phases. Initially, the LLM is prompted to deduce potential causalities or propensities between input features and the described task, utilizing general world knowledge or intuition, and deliberately avoiding reliance on any example cases. This encourages the extraction of relevant prior knowledge related to the problem. In the second phase, the LLM refers to both concrete example demonstrations and insights obtained from the prior phase to formulate class-specific rules. Employing this two-phase structure helps reduce the influence of spurious associations arising from extraneous features and guides the model to prioritize more relevant factors. To balance between underfitting (too few rules) and overfitting (excessively lengthy prompts), we predefine a constant number of rules to be generated for each class (specifically, 10)\footnote{Further analysis on rule count selection appears in Section~\ref{sec:analysis}.}. Moreover, in composing the rules, we instruct the LLM to reference each feature's type as outlined in the dataset's schema, thereby mitigating mismatches of types in the generated rules.

\vspace*{-0.12in}
\paragraph{Response instruction.} To streamline post-processing and application of the generated rules, we incorporate targeted response instructions that guide the LLM in structuring its outputs (see the yellow-highlighted content in Fig.~\ref{fig:prompt_for_rule}). These instructions outline both the format expected for the full response in relation to each output class (i.e., how each answer should be presented) and the specific template that individual rules should conform to (i.e., how to express each rule). With this schema, the LLM selects formatting strategies aligned to the types of features involved. If no further clarification is supplied, the LLM is allowed to join rules using logic operators such as AND or OR. An illustrative example of the output generated with these instructions is provided in Appendix~\ref{sec:outcome-rule}.

\subsection{Inferring Class Likelihood via Rules}
\label{Sec:training}

\paragraph{Parsing rules for feature generation.} In this phase, we employ the rules obtained earlier to construct new binary features for our dataset. For each class, a feature is produced that records whether or not a sample adheres to the class’s associated rule (using a '0' or '1' indicator). These features serve as indicators for the model to estimate class membership probabilities. Nonetheless, given that LLM-derived rules are articulated in natural language, automated feature extraction necessitates an effective parsing methodology. Although we restrict response and rule formats to facilitate parsing during generation, outputs from the LLM may still include inconsistencies or unexpected syntactic variation~\cite{kovavc2023large}; for example, replacing symbols such as ``$>$” or ``$<$” with their verbal equivalents like ``greater than" or ``less than." Such variations can complicate the parsing process.

Rather than resorting to elaborate software engineering solutions to tackle parsing ambiguities, we capitalize on the capabilities of the LLM itself. Because the LLM is adept at comprehending intended meaning despite superficial textual differences and can also synthesize code when provided with explicit guidance (as it is often trained on programming datasets), we format our prompt to include the function’s name, its expected inputs and outputs, and the relevant parsed rules. The prompt is then supplied to the LLM. Illustrative prompts and the results they yield can be found in Figures~\ref{fig:prompt_for_parsing} and~\ref{sec:example_parsing}-\ref{sec:example_parse_outcome} of the Appendix. The LLM-generated code is subsequently run with Python’s \texttt{exec()} to facilitate the conversion of data according to the generated features.

\vspace*{-0.12in}
\paragraph{Inferring class likelihood.} When new binary features are generated for each class based on derived rules, a straightforward approach to estimate the likelihood that a sample belongs to a given class is by tallying the number of rules satisfied per class—that is, summing the respective binary features. Yet, the contribution of each rule and feature may differ in predictive strength, and thus it becomes necessary to quantify their relative importance using the training set. To accomplish this, we introduce a standard machine learning (ML) model tailored to each class’s binary representation. For illustration, we employ a linear model with no intercept term. Let $\mathbf{z}_k^i \in \{0, 1\}^R$ represent the vector of binary features induced from instance $\mathbf{x}^i$ for class $k$, where $R$ is the rule count per class and $c$ denotes the total number of classes. The per-class probability vector, $\mathbf{p}^i$, is formulated as:

\begin{align}
\text{logit}_k^i &= \max(\mathbf{w}_k, 0) \cdot \mathbf{z}_k^i, \nonumber \\
\mathbf{p}^i &= \text{Softmax}({[\text{logit}_{1}^i, ..., \text{logit}_{c}^i]}). \label{eq:infer_prob}
\end{align}

Here, $\mathbf{w}_k$ designates the learnable parameters in the linear model, encapsulating the significance of each feature. By constraining $\mathbf{w}_k$ to be non-negative (i.e., projecting into the positive orthant), we guarantee that the LLM-derived features cannot negatively impact class likelihoods during training. Subsequently, the calculated logit vector is mapped to probability space using the softmax function. To tune $\mathbf{w}_k$, cross-entropy loss is minimized on a limited set of labeled training examples.

\vspace*{-0.12in}
\paragraph{Ensembling with bagging.} To further boost performance, the entire modeling workflow is repeated multiple times, producing an ensemble of inference models. With $T$ runs generating $T$ sets of trained weights, $\{\mathbf{w}[t]\}_{t=1}^T$, classification of a sample is performed by averaging the class probabilities $\mathbf{p}^i$, each derived using one of the weight vectors $\mathbf{w}[t]$ as described in Eq.~\ref{eq:infer_prob}.

To instill stochasticity into rule creation during ensembling, we configure the LLM inference to operate at a sufficiently elevated temperature. Additionally, we permute the sequence of few-shot examples, as empirical findings suggest that the ordering of demonstrations can influence LLM responses~\cite{lu2022fantastically}\footnote{A detailed analysis of rule diversity is provided in Appendix~\ref{sec:diversity}}. To leverage the variability present in predictions and thereby improve robustness, we employ bagging, selecting distinct subsets of features or data instances in each iteration—a strategy that proves particularly advantageous when handling extensive training data or feature sets. This ensemble methodology confers two principal benefits. Firstly, inaccuracies in rule generation by the LLM during certain iterations are mitigated by accurate rules from other rounds, thus strengthening the overall reliability of the framework. This effect aligns with the concept of LLM self-consistency~\cite{wang2022self}: rules that consistently emerge across multiple ensemble members are more probable to be reliable. Secondly, the ensemble mitigates prompt length constraints inherent to LLMs. With input data exceeding the permissible prompt capacity, bagging helps to curtail input size and sustains the model’s robust operation. While an isolated rule set may fail to model the contributions of all features, aggregating multiple weak learners generated across different trials compensates for gaps in feature or instance coverage, neutralizing their impact on individual trials.

\section{Experiments}
We assess the performance of \textsf{FeatLLM} using a suite of tabular datasets and conduct comprehensive analyses to gain insight into the behavior of our approach. Due to limitations of space, extended results, including explorations across various LLM architectures, qualitative studies, and further experimental outcomes, are provided in the Appendix.

\subsection{Performance Evaluation}
\paragraph{Datasets.}
Our experiments draw upon 13 distinct datasets geared toward binary or multi-class classification tasks, including: (1) Adult~\cite{asuncion2007uci}—prediction of annual income exceeding \$50,000; (2) Bank~\cite{moro2014data}—determining the likelihood of a client subscribing to a term deposit; (3) Blood~\cite{yeh2009knowledge}—anticipating repeated donations by blood donors; (4) Car~\cite{kadra2021well}—assessing automobile quality; (5) Communities~\cite{misc_communities_and_crime_183}—forecasting crime rates in different localities; (6) Credit-g~\cite{kadra2021well}—judging individual credit risk as favorable or unfavorable; (7) Diabetes\footnote{\texttt{kaggle.com/datasets/uciml/pima-indians-diabetes-database}}—detecting diabetes in individuals; (8) Heart\footnote{\texttt{kaggle.com/datasets/fedesoriano/heart-failure-prediction}}—predicting the occurrence of coronary artery disease; and (9) Myocardial~\cite{misc_myocardial_infarction_complications_579}—identifying the presence of chronic heart failure.

Our study expands the dataset selection by including several recently released and synthetic datasets that have received limited attention in prior research and were excluded from the pre-training set of the LLMs employed: (10) Cultivars, used to predict the grain yield of specific soybean cultivars\footnote{\texttt{archive.ics.uci.edu/dataset/913}}; (11) NHANES, tasked with identifying the age group of a person based on individual records\footnote{\texttt{archive.ics.uci.edu/dataset/887}}; (12) Sequence-type, which involves classifying numeric sequences into one of four predefined categories: arithmetic, geometric, Fibonacci, or Collatz; (13) Solution-mix, which aims to determine whether a four-solution mixture's concentration percentage surpasses 0.5. The first two datasets have only recently been added to the UCI Machine Learning Repository (after September 2023), guaranteeing their absence from the pre-training stage of the utilized LLM API (gpt-3.5-turbo-0613). The other two datasets are synthetic, designed to preclude any possibility of memorization by the LLM API, and were thus suitable for assessment.

The included datasets exhibit diversity in both their size and complexity, which are further detailed in Table~\ref{Tab:basic-desc}. Each dataset is accompanied by explicit attribute names and comprehensive descriptions. Some datasets, such as `Communities' and `Myocardial,' encompass over a hundred features, challenging the capacity limitations of LLM input windows.

\vspace*{-0.12in}
\paragraph{Baselines.}
We benchmark our approach against nine baseline methods. The initial three are established tabular data modeling techniques: (1) Logistic regression (LogReg), (2) XGBoost~\cite{chen2016xgboost}, and (3) RandomForest~\cite{ho1995random}. Two additional baselines leverage unlabeled datasets: (4) SCARF~\cite{bahri2022scarf} and (5) STUNT~\cite{nam2023stunt}. It is important to recognize that amassing substantial unlabeled data can be impractical in many scenarios. Another recent competitor, (6) TabPFN~\cite{hollmann2022tabpfn}, utilizes a wide range of artificially generated data distributions during pretraining. Finally, three LLM-based methods are evaluated: (7) In-context learning (In-context)~\cite{wei2022emergent}, (8) TABLET~\cite{slack2023tablet}, and (9) TabLLM~\cite{hegselmann2023tabllm}. In-context learning involves inserting a limited set of labeled examples as prompts without adjusting model parameters. TABLET enhances these prompts with supplementary information derived from rule sets and prototype examples through an auxiliary classifier to boost inference accuracy. TabLLM introduces parameter-efficient fine-tuning such as the IA3 technique~\cite{liu2022few}, adapting the LLM to new tasks using a small number of training examples.

\vspace*{-0.12in}
\paragraph{Implementation details.}
\textsf{FeatLLM} is built upon GPT-3.5 as the foundational LLM, but its architecture remains flexible, supporting different LLM choices (refer to Appendix~\ref{sec:backbones} for experiments with alternative backbones). The LLM inference employs a temperature of 0.5 and retains the default API setting for top-p at 1. We configure ensemble size and the count of rules for extraction to be 20 and 10, respectively. An analysis of how hyper-parameters influence results is provided in Figure~\ref{fig:hyperparameter2} and further detailed in Appendix~\ref{sec:hyper-parameter}.
The linear model is trained using the Adam optimizer with a learning rate of 0.01 for a total of 200 epochs, utilizing $k$-fold cross-validation to determine the optimal number of training epochs. For baseline comparisons, GPT-3.5 supports in-context learning, while TabLLM leverages T0~\cite{sanh2022multitask}, consistent with the original experimental setup. It should be noted that TabLLM is limited to LLMs with open-source checkpoints and, consequently, cannot utilize GPT-3.5. With traditional machine learning algorithms like LogReg, XGBoost, and RandomForest, parameter tuning was accomplished through a combination of Grid-search and $k$-fold cross-validation. The value of $k$ is chosen as 2 or 4, ensuring that each class is represented in the training folds. Expanded implementation details can be found in Appendix~\ref{sec:baselines}.

\vspace*{-0.12in}
\paragraph{Results.}
Tables~\ref{tab:main1} and \ref{tab:main2} display a comparative analysis of model performances across multiple datasets. All experiments are conducted three times using different random splits for the training and testing sets, reporting the mean and standard deviation for the AUC metrics. \textsf{FeatLLM} emerges as the leading or second-best method in most evaluations relative to the baseline models. Notably, our framework delivers similar results with only 4 shots, matching those of conventional tabular baselines that demand 32 shots (refer to Fig.~\ref{fig:compares}). While STUNT and TabPFN display marked improvements on particular datasets, their aggregate advantage over traditional machine learning models is limited. Additionally, methods that leverage LLMs, such as in-context learning, TABLET, or TabLLM, are unable to process tabular data with extensive feature dimensions. Through the integration of bagging and ensemble strategies, our approach demonstrates robust performance even on high-dimensional datasets. It is worth mentioning that extending bagging and ensemble methods to alternative LLM-based frameworks is feasible, but this modification would require twice as many inference computations per instance, significantly raising the computational demands and reducing practical viability.

\subsection{Analysis \& Discussion}
\label{sec:analysis}
\paragraph{Ablation study.} To determine the influence of individual model components, we conduct ablation experiments considering the following factors: (1) \textsf{-Tuning}: removing the linear model’s weight tuning phase and instead aggregating binary feature values for logit computation; (2) \textsf{-Ensemble}: excluding the ensemble method; (3) \textsf{-Description}: eliminating the utilization of feature descriptions; and (4) \textsf{-Reasoning}: leaving out the Step-1 procedure in the reasoning instruction phase of rule generation.

The resulting changes in average AUC across all datasets are reported in Table~\ref{tab:ablation}. Across all variants, the removal or modification of each component consistently reduces performance. Remarkably, as the number of shots increases, the advantages of performing weight tuning become more pronounced—suggesting that larger training sets enable more accurate rule importance estimation, thus magnifying the impact of tuning. In contrast, the roles of feature descriptions and reasoning instructions are especially significant for small shot scenarios, highlighting the value of leveraging LLMs’ prior knowledge in low-data regimes. Finally, our ensemble approach delivers steady and noticeable performance gains in every experimental condition.

\vspace*{-0.12in}
\paragraph{Training \& inference time comparison.} We present a comparison of both training and inference durations across several techniques. The evaluation is performed using the Adult dataset, with the training phase utilizing a set of 16 shots. Unless otherwise specified, a single A100 GPU serves as the standard computational resource; TabLLM is the exception, employing four A100 GPUs to enable model parallelism. Table~\ref{tab:cost} details the comparative runtimes. It is worth highlighting that \textsf{FeatLLM} achieves notably efficient inference times, on par with traditional machine learning baselines. In comparison, inference with other LLM-driven frameworks—including In-context learning, TABLET, and TabLLM—demands substantially more time. Regarding training duration, \textsf{FeatLLM} aligns closely with its LLM-based peers, as it necessitates merely 30 API calls, a lightweight demand that does not considerably impact budget and lends itself well to parallel execution.

\vspace*{-0.12in}
\paragraph{Quality of features generated by \textsf{FeatLLM}.}
To assess the utility of the rule-derived features produced by \textsf{FeatLLM} for real-world applications, we run a straightforward experiment. For this, the mean AUROC between each generated feature and the ground-truth target is computed, after which we determine the proportion of generated features that yield an AUROC exceeding 0.5 for each dataset considered. The results, reported in Table~\ref{tab:quality}, indicate that a significant portion of the constructed rules are indeed pertinent to the target variable and enhance task-related insights (see ``Before tuning" column). In addition, when fitting a linear model, the methodology inherently eliminates features with limited predictive value (specifically, those whose learned weights fall below a defined threshold), as indicated in the ``After tuning" column. Here, the threshold has been set to 0.1, based on an analysis of the weight distribution histogram.

\vspace*{-0.12in}
\paragraph{Effect of adding spurious correlations.} To evaluate the ability of different approaches to mitigate the influence of extraneous spurious correlations under low-shot settings, we carried out a targeted experiment. Specifically, we selected the Heart dataset as a base and artificially augmented it by appending randomly chosen columns from the Adult dataset, which are unrelated to the Heart dataset's predictive goal. Model performance was then tracked as we systematically increased the quantity of these irrelevant columns. To ensure comparability across methods, we standardized the experiments to use 8 labeled examples, omitting any methods that require unlabeled data such as SCARF and STUNT.

Figure~\ref{fig:spurious} illustrates the variation in model performance as the extent of spurious correlations increases. Among the models evaluated, \textsf{FeatLLM} demonstrated superior resistance to performance decline, maintaining accuracy as extraneous information increased. This resilience can be attributed to the framework's mechanism of integrating prior knowledge to explicitly align the rule extraction process with relevant task-feature associations, thereby discouraging the creation of rules tied to irrelevant attributes and enhancing robustness. Empirically, rules linked to irrelevant, noisy columns contributed only 1-2\% of the extracted rules. However, we do note that \textsf{FeatLLM} remains susceptible to learning spurious correlations if irrelevant columns stem from datasets with intrinsically similar features, potentially confusing causal inference between task and inputs (see Appendix~\ref{sec:spurious-appendix}).

\vspace*{-0.12in}
\paragraph{Hyper-parameter analysis}
\textsf{FeatLLM} involves two primary hyper-parameters: the ensemble count and the quantity of rules extracted per class in each LLM query. To better understand their effects on model performance, we carried out a series of experiments varying these hyper-parameters. Our observations indicate that increasing the ensemble number yields improved results up to a threshold (approximately 20), after which additional ensembles produce diminishing returns (refer to Fig.~\ref{fig:appendix-hyperparameter1} in Appendix). In examining the number of rules per class, our findings reveal no statistically significant changes in performance across tested values; nevertheless, selecting 10 rules was most advantageous when balancing prompt length with predictive accuracy (refer to Fig.~\ref{fig:hyperparameter2}). We hypothesize that augmenting the number of rules boosts the dimensionality of inputs, which enhances the richness of available information, yet may also foster overfitting, particularly in low-shot contexts.

\section{Conclusion}
This work introduces \textsf{FeatLLM}, establishing a new benchmark for large language model-based few-shot tabular learning. Rather than applying LLMs directly to make individual predictions, \textsf{FeatLLM} leverages them to derive predictive criteria by emulating feature engineering processes. The methodology incorporates rule induction and an ensemble strategy via bagging, enabling highly efficient inference, reliance solely on API interactions (no model retraining), and the ability to address limitations in data feature space size. Through these mechanisms, \textsf{FeatLLM} secures robust predictive performance, marking it as an effective and resource-conscious option for leveraging LLMs on real-world tabular problems.

While \textsf{FeatLLM} currently targets scenarios characterized by limited labeled data (low-shot learning), we intend to further develop the system to automate feature construction in datasets with more ample sample sizes, accommodate diverse feature engineering paradigms beyond rules, and provide greater transparency and interpretability—thereby broadening its practical utility across multiple application domains.

\section{Broader Impact} The \textsf{FeatLLM} framework seeks to exploit the extensive prior knowledge encoded in LLMs to augment tabular learning performance, exceeding the results typically attainable with conventional supervised methods. By enabling more data-efficient learning, our approach contributes toward alleviating overfitting in settings with scarce annotated examples through knowledge transfer from LLMs. This proves especially advantageous for users facing data scarcity and high labeling burdens, as seen in sectors like finance or healthcare, where pretrained LLMs frequently encapsulate field-specific expertise and terminology. It is crucial, however, to recognize the necessity for future work that systematically evaluates how inherent biases or inaccuracies present in LLM knowledge could influence model outputs—especially given scenarios in which model predictions might rely predominantly on such transferred knowledge over limited labeled samples.

\end{document}